% Financial Literacy - Sample Test
%----------------------------------------------
\documentclass[12pt]{exam}

\usepackage{setspace}
\setstretch{1.15}
\nopointsinmargin
\pointformat{}

\usepackage{problemset}

\usepackage[margin=1in]{geometry}

% Uncomment ONE of these:
\noprintanswers
%\printanswers

\begin{document}

% \noindent

% \ExamMode

\FontEasyRead

\Heading
{Financial Literacy}
{Sample Test}
{75 Minutes, 100+ points.
  Up to 5 pages of notes and calculator permitted.
  Answer all multiple-choice questions by selecting the best answer.
  Answer all short-response questions clearly, using correct spelling, punctuation, and grammar.
Show all work for financial math questions.}

\section*{Part I: Multiple Choice (2 points each)}

\begin{questions}

  \begin{mcquestion}[2]
    The tendency to attach more worth to something just because you already own it is called
    \begin{choices}
      \choice Opportunity cost
      \CorrectChoice Endowment effect
      \choice Diversification
      \choice Inflation
    \end{choices}
  \end{mcquestion}

  \begin{mcquestion}[2]
    Gross pay refers to:
    \begin{choices}
      \choice Income after taxes
      \CorrectChoice Income before deductions
      \choice Hourly wage only
      \choice Total yearly salary after bonuses
    \end{choices}
  \end{mcquestion}

  \begin{mcquestion}[2]
    FDIC insurance protects:
    \begin{choices}
      \choice Stock investments
      \CorrectChoice Bank deposits up to legal limits
      \choice Credit card purchases
      \choice Retirement accounts
    \end{choices}
  \end{mcquestion}

  \begin{mcquestion}[2]
    A stock represents:
    \begin{choices}
      \choice A loan to a company
      \CorrectChoice A share of ownership in a company
      \choice A government bond
      \choice A mutual fund
    \end{choices}
  \end{mcquestion}

  \begin{mcquestion}[2]
    Compound interest means:
    \begin{choices}
      \choice Interest is calculated only on the initial principal
      \CorrectChoice Interest is calculated on both the initial principal and accumulated interest
      \choice Interest is paid only at the end of the investment period
      \choice Interest rates are compounded annually
    \end{choices}
  \end{mcquestion}

  \begin{mcquestion}[2]
    Diversification helps investors by:
    \begin{choices}
      \choice Guaranteeing profits
      \CorrectChoice Reducing return volatility
      \choice Avoiding taxes
      \choice Eliminating inflation
    \end{choices}
  \end{mcquestion}

  \begin{mcquestion}[2]
    Why is starting retirement savings early important?
    \begin{choices}
      \choice Taxes are eliminated
      \choice You can always retire earlier
      \CorrectChoice Compound growth has more time to work
      \choice Your employer will cover everything
    \end{choices}
  \end{mcquestion}

  \begin{mcquestion}[2]
    In connection with a car loan or a credit card balance, APR stands for:
    \begin{choices}
      \choice Annual Payment Rate
      \choice Average Principal Return
      \CorrectChoice Annual Percentage Rate
      \choice Applied Payment Rule
    \end{choices}
  \end{mcquestion}

  \begin{mcquestion}[2]
    The main purpose of insurance is to:
    \begin{choices}
      \choice Build wealth
      \CorrectChoice Transfer financial risk
      \choice Avoid taxes
      \choice Raise credit scores
    \end{choices}
  \end{mcquestion}

  \begin{mcquestion}[2]
    Why are credit cards potentially dangerous for many people?
    \begin{choices}
      \choice They cannot be used online
      \choice They are only available for those age 21 and older
      \CorrectChoice High interest debt can accumulate quickly
      \choice They replace debit cards
    \end{choices}
  \end{mcquestion}

  \begin{mcquestion}[2]
    The main purpose of an emergency fund is:
    \begin{choices}
      \choice Vacation savings
      \choice Retirement planning
      \CorrectChoice To cover unexpected expenses
      \choice For investing in stocks
    \end{choices}
  \end{mcquestion}

  \begin{mcquestion}[2]
    Which action best protects consumers from identity theft?
    \begin{choices}
      \choice Using the same password everywhere
      \choice Sharing debit card information with friends
      \CorrectChoice Monitoring financial accounts regularly
      \choice Posting personal information online
    \end{choices}
  \end{mcquestion}

  \begin{mcquestion}[2]
    Which worker would most likely earn overtime pay?
    \begin{choices}
      \choice A salaried executive employee
      \choice A self-employed investor
      \CorrectChoice An hourly employee working more than 40 hours in a week
      \choice A retiree collecting Social Security
    \end{choices}
  \end{mcquestion}

  \begin{mcquestion}[2]
    Which savings vehicle generally offers the highest liquidity?
    \begin{choices}
      \choice Certificate of deposit (CD)
      \CorrectChoice Bank savings account
      \choice 401(k) account
      \choice Stock portfolio
    \end{choices}
  \end{mcquestion}

  \begin{mcquestion}[2]
    Which situation is most likely to trigger an overdraft fee?
    \begin{choices}
      \choice Depositing a paycheck
      \CorrectChoice Writing a check for more than your account balance
      \choice Opening a savings account
      \choice Using direct deposit
    \end{choices}
  \end{mcquestion}

  \begin{mcquestion}[2]
    Which of these investments is generally the most volatile?
    \begin{choices}
      \choice Savings account
      \choice U.S. Treasury bond
      \CorrectChoice Stock in a single company
      \choice Certificate of deposit
    \end{choices}
  \end{mcquestion}

  \begin{mcquestion}[2]
    Which type of insurance would help replace income if someone becomes unable to work because of injury?
    \begin{choices}
      \choice Auto insurance
      \CorrectChoice Disability insurance
      \choice Renters insurance
      \choice Collision insurance
    \end{choices}
  \end{mcquestion}

  \begin{mcquestion}[2]
    Which of these investment accounts often provides tax advantages
    specifically for retirement savings?
    \begin{choices}
      \choice Checking account
      \choice Auto loan
      \CorrectChoice Roth IRA
      \choice Debit card account
    \end{choices}
  \end{mcquestion}

  \begin{mcquestion}[2]
    Which of the following is most likely to \underline{increase} car insurance premiums?
    \begin{choices}
      \choice Maintaining a clean driving record
      \choice Choosing a higher deductible
      \CorrectChoice Having multiple speeding tickets
      \choice Driving fewer miles each year
    \end{choices}
  \end{mcquestion}

  \begin{mcquestion}[2]
    Why is it potentially risky to co-sign a loan for someone else?
    \begin{choices}
      \choice It automatically lowers your taxes
      \choice It hurts your credit score immediately
      \CorrectChoice You become legally responsible if the other person fails to pay
      \choice It eliminates interest charges
    \end{choices}
  \end{mcquestion}

  %%%%%%%%%%%%%%%%%%%%%%%%%%%%%%%%%%%%%%%%%%%%%%%%%%%%%%
  %  \newpage
  \section*{Part II: Short Answer (10 points each)}

  \begin{shortanswer}[10]{2in}
    Taylor earns \$18 per hour and works 25 hours per week.

    Explain:
    \begin{itemize}
      \item Taylor's gross pay for one week
      \item Why gross income differs from take-home pay
      \item Two deductions Taylor likely sees on his paycheck
    \end{itemize}
  \end{shortanswer}
  \begin{solution}
    Gross pay is income before deductions. Take-home pay is lower because taxes and other deductions are removed. Common deductions include federal tax, state tax, Social Security, and Medicare.
  \end{solution}

  \begin{shortanswer}[10]{2in}
    Jordan has \$2,500 saved and is deciding whether to keep it in a checking account, a savings account, or in stocks.
    Explain one advantage and one disadvantage of each savings or investment vehicle.
  \end{shortanswer}
  \begin{solution}
    Checking offers liquidity but low interest. Savings earns interest but limited growth. Stocks offer growth potential but involve risk.
  \end{solution}

  \begin{shortanswer}[10]{2in}
    Two job offers are available:
    \begin{itemize}
      \item Job A: \$19/hour with no benefits
      \item Job B: \$17/hour with health insurance and a 401(k) with a 50\% employer match
    \end{itemize}
    Which job pays better, and why?
  \end{shortanswer}
  \begin{solution}
    Students may discuss benefits, long-term earnings, healthcare costs, retirement contributions, scheduling, or advancement opportunities.
  \end{solution}

  \begin{shortanswer}[10]{3in}
    A student takes out multiple loans for purchases they do not truly need.
    Explain two possible long-term consequences of excessive borrowing.
  \end{shortanswer}
  \begin{solution}
    Possible consequences include debt accumulation, damaged credit, missed payments, reduced financial flexibility, and higher interest costs.
  \end{solution}

  %%%%%%%%%%%%%%%%%%%%%%%%%%%%%%%%%%%%%%%%%%%%%%%%%%%%%%
  \section*{Part III: Financial Math Problems (10 points each)}

  \begin{problem}[10]{2in}
    You deposit \$5,000 into a fund and leave it there for 3 years.  It yields compound
    interest of 2\% the first year,
    3\% the second year,
    and 4\% the third year.  How much will you have at the end of three years?
  \end{problem}

  \begin{problem}[10]{2in}
    You buy a one-year zero coupon \$1,000 bond for \$956.00.  When you redeem the bond
    at maturity for \$1,000, what was the annual interest rate you earned over the year?
  \end{problem}

\end{questions}

\end{document}
